% Raggio dall'aria al materiale della lente a contatto (n = 1,41) con incidenza a 45 gradi:
% riflesso a 45 gradi, rifratto a 30 gradi; fra riflesso e rifratto 105 gradi.
% Esercizio ud7-ottica-riflesso-rifratto (soluzione)
\begin{tikzpicture}[>=Stealth, font=\small]
  \fill[formula!10] (-3,-2.7) rectangle (3,0);
  \draw[thick] (-3,0) -- (3,0);
  \node[right] at (3,0.35) {aria};
  \node[right] at (3,-0.4) {$n = 1{,}41$};
  \draw[dashed, gray] (0,2.6) -- (0,-2.7) node[below] {normale};
  % raggi
  \draw[esempio, very thick, ->] (-2.1,2.1) -- (-1.05,1.05);
  \draw[esempio, very thick] (-1.05,1.05) -- (0,0);
  \draw[esempio, very thick, ->] (0,0) -- (1.05,1.05);
  \draw[esempio, very thick] (1.05,1.05) -- (2.1,2.1);
  \draw[esempio, very thick, ->] (0,0) -- (0.65,-1.13);
  \draw[esempio, very thick] (0.65,-1.13) -- (1.3,-2.25);
  \node[esempio, above left] at (-2.1,2.1) {incidente};
  \node[esempio, above right] at (2.1,2.1) {riflesso};
  \node[esempio, right] at (1.3,-2.25) {rifratto};
  % angoli con la normale
  \draw (0,1.0) arc[start angle=90, end angle=135, radius=1.0];
  \node at (112.5:1.4) {\ang{45}};
  \draw (0,1.0) arc[start angle=90, end angle=45, radius=1.0];
  \node at (67.5:1.4) {\ang{45}};
  \draw (0,-1.0) arc[start angle=270, end angle=300, radius=1.0];
  \node at (285:1.45) {\ang{30}};
  % angolo fra riflesso e rifratto
  \draw[attenzione, thick] (45:0.55) arc[start angle=45, end angle=-60, radius=0.55];
  \node[attenzione, anchor=west] at (14:0.6) {\ang{105}};
\end{tikzpicture}
